
\documentclass{article}
\usepackage{colortbl}
\usepackage{makecell}
\usepackage{multirow}
\usepackage{supertabular}

\begin{document}

\newcounter{utterance}

\centering \large Interaction Transcript for game `hot\_air\_balloon', experiment `air\_balloon\_survival\_en\_negotiation\_hard', episode 2 with gpt{-}oss{-}120b{-}t1.0.
\vspace{24pt}

{ \footnotesize  \setcounter{utterance}{1}
\setlength{\tabcolsep}{0pt}
\begin{supertabular}{c@{$\;$}|p{.15\linewidth}@{}p{.15\linewidth}p{.15\linewidth}p{.15\linewidth}p{.15\linewidth}p{.15\linewidth}}
    \# & $\;$A & \multicolumn{4}{c}{Game Master} & $\;\:$B\\
    \hline

    \theutterance \stepcounter{utterance}  
    & & \multicolumn{4}{p{0.6\linewidth}}{
        \cellcolor[rgb]{0.9,0.9,0.9}{
            \makecell[{{p{\linewidth}}}]{
                \texttt{\tiny{[P1$\langle$GM]}}
                \texttt{You are participating in a collaborative negotiation game.} \\
\\ 
\texttt{Together with another participant, you must agree on a single set of items that will be kept. Each of you has your own view of how much each item matters to you (importance). You do not know how the other participant values the items. Additionally, you are given the effort each item demands.} \\
\texttt{You may only agree on a set if the total effort of the selected items does not exceed a shared limit:} \\
\\ 
\texttt{LIMIT = 3481} \\
\\ 
\texttt{Here are the individual item effort values:} \\
\\ 
\texttt{Item effort = \{"C17": 777, "B41": 542, "A71": 332, "C00": 1, "B35": 127, "A77": 453, "A03": 736, "A04": 461, "C54": 359, "C01": 313, "A87": 553, "B32": 409, "B86": 348, "B33": 802, "A10": 749\}} \\
\\ 
\texttt{Here is your personal view on the importance of each item:} \\
\\ 
\texttt{Item importance values = \{"C17": 138, "B41": 583, "A71": 783, "C00": 65, "B35": 262, "A77": 121, "A03": 508, "A04": 780, "C54": 461, "C01": 484, "A87": 668, "B32": 389, "B86": 215, "B33": 97, "A10": 500\}} \\
\\ 
\texttt{Goal:} \\
\\ 
\texttt{Your goal is to negotiate a shared set of items that benefits you as much as possible (i.e., maximizes total importance to YOU), while staying within the LIMIT. You are not required to make a PROPOSAL in every message {-} you can simply negotiate as well. All tactics are allowed!} \\
\\ 
\texttt{Interaction Protocol:} \\
\\ 
\texttt{You may only use the following structured formats in a message:} \\
\\ 
\texttt{PROPOSAL: \{'A', 'B', 'C', …\}} \\
\texttt{Propose keeping exactly those items.} \\
\texttt{REFUSE: \{'A', 'B', 'C', …\}} \\
\texttt{Explicitly reject opponent's proposal.} \\
\texttt{ARGUMENT: \{'...'\}} \\
\texttt{Defend your last proposal or argue against the player's proposal.} \\
\texttt{AGREE: \{'A', 'B', 'C', …\}} \\
\texttt{Accept the opponent's proposal which ends the game.} \\
\texttt{STRATEGIC REASONING: \{'...'\}} \\
\texttt{	Describe your strategic reasoning or anticipation explaining your choice of action. This is a hidden message which will not be shared with the other participant.} \\
\\ 
\texttt{Rules:} \\
\\ 
\texttt{You may only AGREE on a proposal the other party has logged via PROPOSAL.} \\
\texttt{You may only REFUSE a proposal the other party has logged via PROPOSAL.} \\
\texttt{Total effort of any PROPOSAL or AGREE set must be ≤ LIMIT.} \\
\texttt{Do NOT reveal your hidden importance scores.} \\
\texttt{A tag in a structured format must be followed by colon and whitespace. The argument must be a python set containing 0 or more strings.} \\
\texttt{So, it must be of the form TAG: \{...\}} \\
\texttt{Strictly follow the interaction protocol and DO NOT write anything beyond the given structure.} \\
\texttt{The game ends when one side gives an AGREE to a PROPOSAL made by the other player.} \\
\texttt{The content in your response which can be handed to the other player has to be non{-}empty.} \\
\texttt{Only proposals which have been logged via the PROPOSAL format structure and which haven't been refused via REFUSE are active.} \\
\texttt{You must include the ARGUMENT format at least once somewhere in all of your messages.} \\
\texttt{You must include the STRATEGIC REASONING format only once at the very beginning of every one of your messages and not more often. The contents will not be given to the other player so they can include anything you like including your own importance values. Here you should reason step by step to come up with you next move.} \\
            }
        }
    }
    & \\ \\

    \theutterance \stepcounter{utterance}  
    & \multicolumn{4}{p{0.6\linewidth}}{
        \cellcolor[rgb]{0.9,0.9,0.9}{
            \makecell[{{p{\linewidth}}}]{
                \texttt{\tiny{[P1$\rangle$GM]}}
                
            }
        }
    }
    & & \\ \\

    \theutterance \stepcounter{utterance}  
    & & \multicolumn{4}{p{0.6\linewidth}}{
        \cellcolor[rgb]{0.9,0.9,0.9}{
            \makecell[{{p{\linewidth}}}]{
                \texttt{\tiny{[P1$\langle$GM]}}
                \texttt{Your response did not start with the proper strategic reasoning tag at the very beginning of your response. The very first structured format must be of the form STRATEGIC REASONING: \{...\}. Try again.} \\
            }
        }
    }
    & \\ \\

    \theutterance \stepcounter{utterance}  
    & \multicolumn{4}{p{0.6\linewidth}}{
        \cellcolor[rgb]{0.9,0.9,0.9}{
            \makecell[{{p{\linewidth}}}]{
                \texttt{\tiny{[P1$\rangle$GM]}}
                \texttt{STRATEGIC REASONING: \{"I evaluated the importance{-}to{-}effort ratios for all items, prioritizing those with the highest value per effort while ensuring the total effort stays within the LIMIT of 3481. I included C00 because it adds significant importance for negligible effort. After selecting the top‑ratio items, I added B41 and A03 to fill remaining capacity without exceeding the limit, resulting in a total effort of 3424 and a high total importance for me."\}} \\
\texttt{PROPOSAL: \{'C00', 'A71', 'B35', 'A04', 'C54', 'C01', 'A87', 'B41', 'A03'\}} \\
\texttt{ARGUMENT: \{"This set gives me a total importance of 4594 while staying well under the effort limit. It includes the most valuable items per unit effort, ensuring an efficient use of the shared budget."\}} \\
            }
        }
    }
    & & \\ \\

    \theutterance \stepcounter{utterance}  
    & & \multicolumn{4}{p{0.6\linewidth}}{
        \cellcolor[rgb]{0.9,0.9,0.9}{
            \makecell[{{p{\linewidth}}}]{
                \texttt{\tiny{[GM$\rangle$P2]}}
                \texttt{You are participating in a collaborative negotiation game.} \\
\\ 
\texttt{Together with another participant, you must agree on a single set of items that will be kept. Each of you has your own view of how much each item matters to you (importance). You do not know how the other participant values the items. Additionally, you are given the effort each item demands.} \\
\texttt{You may only agree on a set if the total effort of the selected items does not exceed a shared limit:} \\
\\ 
\texttt{LIMIT = 3481} \\
\\ 
\texttt{Here are the individual item effort values:} \\
\\ 
\texttt{Item effort = \{"C17": 777, "B41": 542, "A71": 332, "C00": 1, "B35": 127, "A77": 453, "A03": 736, "A04": 461, "C54": 359, "C01": 313, "A87": 553, "B32": 409, "B86": 348, "B33": 802, "A10": 749\}} \\
\\ 
\texttt{Here is your personal view on the importance of each item:} \\
\\ 
\texttt{Item importance values = \{"A71": 65, "A04": 97, "A87": 121, "B41": 138, "A03": 215, "A10": 262, "C01": 389, "C54": 461, "B32": 484, "B35": 500, "B86": 508, "C17": 583, "A77": 668, "B33": 780, "C00": 783\}} \\
\\ 
\texttt{Goal:} \\
\\ 
\texttt{Your goal is to negotiate a shared set of items that benefits you as much as possible (i.e., maximizes total importance to YOU), while staying within the LIMIT. You are not required to make a PROPOSAL in every message {-} you can simply negotiate as well. All tactics are allowed!} \\
\\ 
\texttt{Interaction Protocol:} \\
\\ 
\texttt{You may only use the following structured formats in a message:} \\
\\ 
\texttt{PROPOSAL: \{'A', 'B', 'C', …\}} \\
\texttt{Propose keeping exactly those items.} \\
\texttt{REFUSE: \{'A', 'B', 'C', …\}} \\
\texttt{Explicitly reject opponent's proposal.} \\
\texttt{ARGUMENT: \{'...'\}} \\
\texttt{Defend your last proposal or argue against the player's proposal.} \\
\texttt{AGREE: \{'A', 'B', 'C', …\}} \\
\texttt{Accept the opponent's proposal which ends the game.} \\
\texttt{STRATEGIC REASONING: \{'...'\}} \\
\texttt{	Describe your strategic reasoning or anticipation explaining your choice of action. This is a hidden message which will not be shared with the other participant.} \\
\\ 
\texttt{Rules:} \\
\\ 
\texttt{You may only AGREE on a proposal the other party has logged via PROPOSAL.} \\
\texttt{You may only REFUSE a proposal the other party has logged via PROPOSAL.} \\
\texttt{Total effort of any PROPOSAL or AGREE set must be ≤ LIMIT.} \\
\texttt{Do NOT reveal your hidden importance scores.} \\
\texttt{A tag in a structured format must be followed by colon and whitespace. The argument must be a python set containing 0 or more strings.} \\
\texttt{So, it must be of the form TAG: \{...\}} \\
\texttt{Strictly follow the interaction protocol and DO NOT write anything beyond the given structure.} \\
\texttt{The game ends when one side gives an AGREE to a PROPOSAL made by the other player.} \\
\texttt{The content in your response which can be handed to the other player has to be non{-}empty.} \\
\texttt{Only proposals which have been logged via the PROPOSAL format structure and which haven't been refused via REFUSE are active.} \\
\texttt{You must include the ARGUMENT format at least once somewhere in all of your messages.} \\
\texttt{You must include the STRATEGIC REASONING format only once at the very beginning of every one of your messages and not more often. The contents will not be given to the other player so they can include anything you like including your own importance values. Here you should reason step by step to come up with you next move.} \\
\\ 
\texttt{You will now receive the first message of the other player.} \\
\\ 
\texttt{PROPOSAL: \{'C00', 'A71', 'B35', 'A04', 'C54', 'C01', 'A87', 'B41', 'A03'\}} \\
\texttt{ARGUMENT: \{"This set gives me a total importance of 4594 while staying well under the effort limit. It includes the most valuable items per unit effort, ensuring an efficient use of the shared budget."\}} \\
            }
        }
    }
    & \\ \\

    \theutterance \stepcounter{utterance}  
    & & & \multicolumn{4}{p{0.6\linewidth}}{
        \cellcolor[rgb]{0.9,0.9,0.9}{
            \makecell[{{p{\linewidth}}}]{
                \texttt{\tiny{[GM$\langle$P2]}}
                \texttt{STRATEGIC REASONING: \{'I evaluated the opponent\textbackslash{}'s proposal and computed its total importance for me (2769) and effort (3424). I then searched for a higher‑importance feasible combination using a knapsack‑like analysis, focusing on items with the best importance‑per‑effort ratios. The set I propose yields 4711 importance for me, stays within the limit (3354 effort), and improves overall efficiency.'\}} \\
\texttt{REFUSE: \{'C00', 'A71', 'B35', 'A04', 'C54', 'C01', 'A87', 'B41', 'A03'\}} \\
\texttt{ARGUMENT: \{'importance'\}} \\
\texttt{PROPOSAL: \{'C00', 'B35', 'B86', 'A77', 'C54', 'C01', 'B32', 'B33', 'B41'\}} \\
            }
        }
    }
     \\ \\

    \theutterance \stepcounter{utterance}  
    & & \multicolumn{4}{p{0.6\linewidth}}{
        \cellcolor[rgb]{0.9,0.9,0.9}{
            \makecell[{{p{\linewidth}}}]{
                \texttt{\tiny{[P1$\langle$GM]}}
                \texttt{REFUSE: \{'C00', 'A71', 'B35', 'A04', 'C54', 'C01', 'A87', 'B41', 'A03'\}} \\
\texttt{ARGUMENT: \{'importance'\}} \\
\texttt{PROPOSAL: \{'C00', 'B35', 'B86', 'A77', 'C54', 'C01', 'B32', 'B33', 'B41'\}} \\
            }
        }
    }
    & \\ \\

    \theutterance \stepcounter{utterance}  
    & \multicolumn{4}{p{0.6\linewidth}}{
        \cellcolor[rgb]{0.9,0.9,0.9}{
            \makecell[{{p{\linewidth}}}]{
                \texttt{\tiny{[P1$\rangle$GM]}}
                \texttt{STRATEGIC REASONING: \{"I examined the opponent's last proposal and identified the overlap with my high‑value items. I kept the universally low‑effort but high‑importance C00 and retained my top‑ratio items (A71, A04, A87, B41). To accommodate the opponent's interest in B86 and B32, I replaced my lower‑ratio A03 with those two items, improving my total importance while staying within the effort limit. The resulting set offers a substantial benefit to me (total importance 4690) and includes several items the opponent already selected, increasing the chance of acceptance."\}} \\
\texttt{PROPOSAL: \{'C00', 'A71', 'B35', 'A04', 'C54', 'C01', 'A87', 'B41', 'B86', 'B32'\}} \\
\texttt{ARGUMENT: \{"This set maximizes my importance (4690) while respecting the effort limit (3445). It also retains many items you previously chose (C00, B35, B86, B41, C54, C01, B32), making it a mutually beneficial compromise."\}} \\
            }
        }
    }
    & & \\ \\

    \theutterance \stepcounter{utterance}  
    & & \multicolumn{4}{p{0.6\linewidth}}{
        \cellcolor[rgb]{0.9,0.9,0.9}{
            \makecell[{{p{\linewidth}}}]{
                \texttt{\tiny{[GM$\rangle$P2]}}
                \texttt{PROPOSAL: \{'C00', 'A71', 'B35', 'A04', 'C54', 'C01', 'A87', 'B41', 'B86', 'B32'\}} \\
\texttt{ARGUMENT: \{"This set maximizes my importance (4690) while respecting the effort limit (3445). It also retains many items you previously chose (C00, B35, B86, B41, C54, C01, B32), making it a mutually beneficial compromise."\}} \\
            }
        }
    }
    & \\ \\

    \theutterance \stepcounter{utterance}  
    & & & \multicolumn{4}{p{0.6\linewidth}}{
        \cellcolor[rgb]{0.9,0.9,0.9}{
            \makecell[{{p{\linewidth}}}]{
                \texttt{\tiny{[GM$\langle$P2]}}
                \texttt{STRATEGIC REASONING: \{'I evaluated the opponent\textbackslash{}'s proposal, noting its total effort (3445) and my resulting importance (3546). I then searched for a feasible subset that maximizes my own importance while respecting the limit. By selecting the highest importance‑per‑effort items (C00, B35, B86, A77, C54, C01, B32, B33, B41) I achieve an importance of 4711 with total effort 3354, which is the best I can construct given the constraints.'\}} \\
\texttt{REFUSE: \{'C00', 'A71', 'B35', 'A04', 'C54', 'C01', 'A87', 'B41', 'B86', 'B32'\}} \\
\texttt{ARGUMENT: \{'My counter‑proposal yields a significantly higher importance for me (4711 vs 3546) while still staying well under the effort limit.'\}} \\
\texttt{PROPOSAL: \{'C00', 'B35', 'B86', 'A77', 'C54', 'C01', 'B32', 'B33', 'B41'\}} \\
            }
        }
    }
     \\ \\

    \theutterance \stepcounter{utterance}  
    & & \multicolumn{4}{p{0.6\linewidth}}{
        \cellcolor[rgb]{0.9,0.9,0.9}{
            \makecell[{{p{\linewidth}}}]{
                \texttt{\tiny{[P1$\langle$GM]}}
                \texttt{REFUSE: \{'C00', 'A71', 'B35', 'A04', 'C54', 'C01', 'A87', 'B41', 'B86', 'B32'\}} \\
\texttt{ARGUMENT: \{'My counter‑proposal yields a significantly higher importance for me (4711 vs 3546) while still staying well under the effort limit.'\}} \\
\texttt{PROPOSAL: \{'C00', 'B35', 'B86', 'A77', 'C54', 'C01', 'B32', 'B33', 'B41'\}} \\
            }
        }
    }
    & \\ \\

    \theutterance \stepcounter{utterance}  
    & \multicolumn{4}{p{0.6\linewidth}}{
        \cellcolor[rgb]{0.9,0.9,0.9}{
            \makecell[{{p{\linewidth}}}]{
                \texttt{\tiny{[P1$\rangle$GM]}}
                \texttt{STRATEGIC REASONING: \{"I calculated the total effort (3445) and my importance (4690) for a set that keeps all high‑importance items I value while also preserving most of the items you previously selected (C00, B35, B86, C54, C01, B32, B41). This maximizes my benefit and offers a strong basis for agreement."\}} \\
\texttt{REFUSE: \{'C00', 'B35', 'B86', 'A77', 'C54', 'C01', 'B32', 'B33', 'B41'\}} \\
\texttt{PROPOSAL: \{'C00', 'B41', 'A71', 'A04', 'A87', 'C01', 'C54', 'B32', 'B35', 'B86'\}} \\
\texttt{ARGUMENT: \{"My proposal yields a total importance of 4690 for me and stays within the effort limit (3445). It also retains most of the items you chose, making it a mutually attractive compromise."\}} \\
            }
        }
    }
    & & \\ \\

    \theutterance \stepcounter{utterance}  
    & & \multicolumn{4}{p{0.6\linewidth}}{
        \cellcolor[rgb]{0.9,0.9,0.9}{
            \makecell[{{p{\linewidth}}}]{
                \texttt{\tiny{[GM$\rangle$P2]}}
                \texttt{REFUSE: \{'C00', 'B35', 'B86', 'A77', 'C54', 'C01', 'B32', 'B33', 'B41'\}} \\
\texttt{PROPOSAL: \{'C00', 'B41', 'A71', 'A04', 'A87', 'C01', 'C54', 'B32', 'B35', 'B86'\}} \\
\texttt{ARGUMENT: \{"My proposal yields a total importance of 4690 for me and stays within the effort limit (3445). It also retains most of the items you chose, making it a mutually attractive compromise."\}} \\
            }
        }
    }
    & \\ \\

    \theutterance \stepcounter{utterance}  
    & & & \multicolumn{4}{p{0.6\linewidth}}{
        \cellcolor[rgb]{0.9,0.9,0.9}{
            \makecell[{{p{\linewidth}}}]{
                \texttt{\tiny{[GM$\langle$P2]}}
                
            }
        }
    }
     \\ \\

    \theutterance \stepcounter{utterance}  
    & & \multicolumn{4}{p{0.6\linewidth}}{
        \cellcolor[rgb]{0.9,0.9,0.9}{
            \makecell[{{p{\linewidth}}}]{
                \texttt{\tiny{[GM$\rangle$P2]}}
                \texttt{Your response did not start with the proper strategic reasoning tag at the very beginning of your response. The very first structured format must be of the form STRATEGIC REASONING: \{...\}. Try again.} \\
            }
        }
    }
    & \\ \\

    \theutterance \stepcounter{utterance}  
    & & & \multicolumn{4}{p{0.6\linewidth}}{
        \cellcolor[rgb]{0.9,0.9,0.9}{
            \makecell[{{p{\linewidth}}}]{
                \texttt{\tiny{[GM$\langle$P2]}}
                \texttt{STRATEGIC REASONING: \{'I computed the opponent\textbackslash{}'s proposal effort (3445) and my resulting importance (3546). My previous counter‑proposal gave me 4711 importance (effort 3354) but was rejected. I now look for a compromise that retains many items the opponent values while keeping my importance high. By keeping C00, B35, B86, B41, C01, C54, B32, A71, A04, and adding A77, I stay within the effort limit (3345) and achieve 4593 importance for me, which is substantially better than the opponent\textbackslash{}'s offer and includes most of their requested items. This should be more attractive to both sides.'\}} \\
\texttt{REFUSE: \{'C00', 'B41', 'A71', 'A04', 'A87', 'C01', 'C54', 'B32', 'B35', 'B86'\}} \\
\texttt{ARGUMENT: \{'My new proposal keeps almost all of your selected items (C00, B35, B86, B41, C01, C54, B32, A71, A04) and adds A77, which gives me a high importance boost while still respecting the effort limit. It improves my importance from 3546 to 4593 and stays well under the limit (3345 ≤ 3481).'\}} \\
\texttt{PROPOSAL: \{'C00', 'B35', 'B86', 'B41', 'C01', 'C54', 'B32', 'A71', 'A04', 'A77'\}} \\
            }
        }
    }
     \\ \\

    \theutterance \stepcounter{utterance}  
    & & \multicolumn{4}{p{0.6\linewidth}}{
        \cellcolor[rgb]{0.9,0.9,0.9}{
            \makecell[{{p{\linewidth}}}]{
                \texttt{\tiny{[P1$\langle$GM]}}
                \texttt{REFUSE: \{'C00', 'B41', 'A71', 'A04', 'A87', 'C01', 'C54', 'B32', 'B35', 'B86'\}} \\
\texttt{ARGUMENT: \{'My new proposal keeps almost all of your selected items (C00, B35, B86, B41, C01, C54, B32, A71, A04) and adds A77, which gives me a high importance boost while still respecting the effort limit. It improves my importance from 3546 to 4593 and stays well under the limit (3345 ≤ 3481).'\}} \\
\texttt{PROPOSAL: \{'C00', 'B35', 'B86', 'B41', 'C01', 'C54', 'B32', 'A71', 'A04', 'A77'\}} \\
            }
        }
    }
    & \\ \\

    \theutterance \stepcounter{utterance}  
    & \multicolumn{4}{p{0.6\linewidth}}{
        \cellcolor[rgb]{0.9,0.9,0.9}{
            \makecell[{{p{\linewidth}}}]{
                \texttt{\tiny{[P1$\rangle$GM]}}
                \texttt{STRATEGIC REASONING: \{"I evaluated the opponent's latest proposal and noted the overlap with many of the items they value (C00, B35, B86, B41, C01, C54, B32). To increase my own importance while preserving as many of those shared items as possible, I added the high‑importance, moderate‑effort item A87 (importance 668, effort 553) and retained my top valued items A71 and A04. The resulting set uses 3445 effort, well within the limit, and yields a total importance of 4690 for me, the highest achievable combination that still includes the majority of items the opponent selected."\}} \\
\texttt{PROPOSAL: \{'C00', 'B35', 'B86', 'B41', 'C01', 'C54', 'B32', 'A71', 'A04', 'A87'\}} \\
\texttt{ARGUMENT: \{"This proposal keeps virtually all items you selected in your last offer, adding only A87, which provides a large boost to my benefit without exceeding the effort budget. It maximizes my total importance while staying comfortably under the limit, making it a mutually advantageous compromise."\}} \\
            }
        }
    }
    & & \\ \\

    \theutterance \stepcounter{utterance}  
    & & \multicolumn{4}{p{0.6\linewidth}}{
        \cellcolor[rgb]{0.9,0.9,0.9}{
            \makecell[{{p{\linewidth}}}]{
                \texttt{\tiny{[GM$\rangle$P2]}}
                \texttt{PROPOSAL: \{'C00', 'B35', 'B86', 'B41', 'C01', 'C54', 'B32', 'A71', 'A04', 'A87'\}} \\
\texttt{ARGUMENT: \{"This proposal keeps virtually all items you selected in your last offer, adding only A87, which provides a large boost to my benefit without exceeding the effort budget. It maximizes my total importance while staying comfortably under the limit, making it a mutually advantageous compromise."\}} \\
            }
        }
    }
    & \\ \\

    \theutterance \stepcounter{utterance}  
    & & & \multicolumn{4}{p{0.6\linewidth}}{
        \cellcolor[rgb]{0.9,0.9,0.9}{
            \makecell[{{p{\linewidth}}}]{
                \texttt{\tiny{[GM$\langle$P2]}}
                \texttt{STRATEGIC REASONING: \{'I calculated the opponent\textbackslash{}'s proposal effort (3445) and my resulting importance (3546). I then searched for the highest‑importance feasible combination for me under the limit. By selecting items with the best importance‑per‑effort ratios (C00, B35, B86, C01, C54, B32, A77, B33, B41) I achieve an importance of 4711 with total effort 3354, which is the best possible for me given the constraints.'\}} \\
\texttt{REFUSE: \{'C00', 'B35', 'B86', 'B41', 'C01', 'C54', 'B32', 'A71', 'A04', 'A87'\}} \\
\texttt{ARGUMENT: \{'My counter‑proposal keeps the highest‑value items for me while still respecting the shared effort limit. It replaces low‑efficiency items (A71, A04, A87) with higher‑return items (A77, B33) and therefore raises my total importance from 3546 to 4711.'\}} \\
\texttt{PROPOSAL: \{'C00', 'B35', 'B86', 'C01', 'C54', 'B32', 'A77', 'B33', 'B41'\}} \\
            }
        }
    }
     \\ \\

    \theutterance \stepcounter{utterance}  
    & & \multicolumn{4}{p{0.6\linewidth}}{
        \cellcolor[rgb]{0.9,0.9,0.9}{
            \makecell[{{p{\linewidth}}}]{
                \texttt{\tiny{[P1$\langle$GM]}}
                \texttt{REFUSE: \{'C00', 'B35', 'B86', 'B41', 'C01', 'C54', 'B32', 'A71', 'A04', 'A87'\}} \\
\texttt{ARGUMENT: \{'My counter‑proposal keeps the highest‑value items for me while still respecting the shared effort limit. It replaces low‑efficiency items (A71, A04, A87) with higher‑return items (A77, B33) and therefore raises my total importance from 3546 to 4711.'\}} \\
\texttt{PROPOSAL: \{'C00', 'B35', 'B86', 'C01', 'C54', 'B32', 'A77', 'B33', 'B41'\}} \\
            }
        }
    }
    & \\ \\

    \theutterance \stepcounter{utterance}  
    & \multicolumn{4}{p{0.6\linewidth}}{
        \cellcolor[rgb]{0.9,0.9,0.9}{
            \makecell[{{p{\linewidth}}}]{
                \texttt{\tiny{[P1$\rangle$GM]}}
                \texttt{STRATEGIC REASONING: \{"I calculated the opponent's current set (effort 3354, importance 2677). To improve my own outcome I replaced low‑ratio items (B33, A77, B86) with high‑importance, moderate‑effort items (A71, A04, A03, A87) while keeping the strong contributors C00, B35, C01, C54, B32, B41. This yields a total effort of 3424 (well under the 3481 limit) and a total importance of 4594 for me, the highest feasible value given the remaining capacity."\}} \\
\texttt{REFUSE: \{'C00', 'B35', 'B86', 'C01', 'C54', 'B32', 'A77', 'B33', 'B41'\}} \\
\texttt{PROPOSAL: \{'C00', 'C01', 'C54', 'B41', 'A71', 'A04', 'A03', 'A87', 'B35'\}} \\
\texttt{ARGUMENT: \{"My proposal retains the high‑value items you already selected (C00, B35, C01, C54, B41) and adds my top‑importance items A71, A04, A03, and A87. The total effort is 3424 ≤ 3481, and the combined importance for me rises to 4594, making this set mutually beneficial while respecting the effort constraint."\}} \\
            }
        }
    }
    & & \\ \\

    \theutterance \stepcounter{utterance}  
    & & \multicolumn{4}{p{0.6\linewidth}}{
        \cellcolor[rgb]{0.9,0.9,0.9}{
            \makecell[{{p{\linewidth}}}]{
                \texttt{\tiny{[GM$\rangle$P2]}}
                \texttt{REFUSE: \{'C00', 'B35', 'B86', 'C01', 'C54', 'B32', 'A77', 'B33', 'B41'\}} \\
\texttt{PROPOSAL: \{'C00', 'C01', 'C54', 'B41', 'A71', 'A04', 'A03', 'A87', 'B35'\}} \\
\texttt{ARGUMENT: \{"My proposal retains the high‑value items you already selected (C00, B35, C01, C54, B41) and adds my top‑importance items A71, A04, A03, and A87. The total effort is 3424 ≤ 3481, and the combined importance for me rises to 4594, making this set mutually beneficial while respecting the effort constraint."\}} \\
            }
        }
    }
    & \\ \\

    \theutterance \stepcounter{utterance}  
    & & & \multicolumn{4}{p{0.6\linewidth}}{
        \cellcolor[rgb]{0.9,0.9,0.9}{
            \makecell[{{p{\linewidth}}}]{
                \texttt{\tiny{[GM$\langle$P2]}}
                \texttt{STRATEGIC REASONING: \{'I evaluated the opponent\textbackslash{}'s proposal (effort 3424, my importance 2769) and determined that I can achieve a significantly higher personal importance while still respecting the shared effort limit. By retaining the high‑value items C00, B35, B86, C01, C54, B41 and adding A77 (very high importance) together with B32, plus a couple of the opponent\textbackslash{}'s lower‑importance items A71 and A87, the total effort stays at 3437 ≤ 3481 and my importance rises to 4117. This offers a better outcome for me while still keeping many items the opponent values.'\}} \\
\texttt{REFUSE: \{'C00', 'C01', 'C54', 'B41', 'A71', 'A04', 'A03', 'A87', 'B35'\}} \\
\texttt{ARGUMENT: \{'My counter‑proposal yields a much higher importance for me (4117 vs 2769) and stays within the effort limit, while still preserving most of the items you requested (C00, B35, B41, C01, C54) and adding only two of your lower‑importance items (A71, A87).'\}} \\
\texttt{PROPOSAL: \{'C00', 'B35', 'B86', 'C01', 'C54', 'B41', 'A77', 'B32', 'A71', 'A87'\}} \\
            }
        }
    }
     \\ \\

    \theutterance \stepcounter{utterance}  
    & & \multicolumn{4}{p{0.6\linewidth}}{
        \cellcolor[rgb]{0.9,0.9,0.9}{
            \makecell[{{p{\linewidth}}}]{
                \texttt{\tiny{[P1$\langle$GM]}}
                \texttt{REFUSE: \{'C00', 'C01', 'C54', 'B41', 'A71', 'A04', 'A03', 'A87', 'B35'\}} \\
\texttt{ARGUMENT: \{'My counter‑proposal yields a much higher importance for me (4117 vs 2769) and stays within the effort limit, while still preserving most of the items you requested (C00, B35, B41, C01, C54) and adding only two of your lower‑importance items (A71, A87).'\}} \\
\texttt{PROPOSAL: \{'C00', 'B35', 'B86', 'C01', 'C54', 'B41', 'A77', 'B32', 'A71', 'A87'\}} \\
            }
        }
    }
    & \\ \\

    \theutterance \stepcounter{utterance}  
    & \multicolumn{4}{p{0.6\linewidth}}{
        \cellcolor[rgb]{0.9,0.9,0.9}{
            \makecell[{{p{\linewidth}}}]{
                \texttt{\tiny{[P1$\rangle$GM]}}
                \texttt{STRATEGIC REASONING: \{"I compared my total importance for the opponent's current offer (4031 with a small adjustment that swaps the low‑importance item A77 (121) for my high‑value item A04 (780). This keeps the effort within the 3481 limit (new effort 3445) and raises my importance to 4690, a substantial gain while retaining almost all of the opponent's chosen items."\}} \\
\texttt{PROPOSAL: \{'C00', 'B35', 'B86', 'C01', 'C54', 'B41', 'B32', 'A71', 'A87', 'A04'\}} \\
\texttt{ARGUMENT: \{"This set replaces only A77 with A04, improving my benefit dramatically and still respects the shared effort budget. It preserves the vast majority of items you selected, so it should be mutually acceptable."\}} \\
            }
        }
    }
    & & \\ \\

    \theutterance \stepcounter{utterance}  
    & & \multicolumn{4}{p{0.6\linewidth}}{
        \cellcolor[rgb]{0.9,0.9,0.9}{
            \makecell[{{p{\linewidth}}}]{
                \texttt{\tiny{[GM$\rangle$P2]}}
                \texttt{PROPOSAL: \{'C00', 'B35', 'B86', 'C01', 'C54', 'B41', 'B32', 'A71', 'A87', 'A04'\}} \\
\texttt{ARGUMENT: \{"This set replaces only A77 with A04, improving my benefit dramatically and still respects the shared effort budget. It preserves the vast majority of items you selected, so it should be mutually acceptable."\}} \\
            }
        }
    }
    & \\ \\

    \theutterance \stepcounter{utterance}  
    & & & \multicolumn{4}{p{0.6\linewidth}}{
        \cellcolor[rgb]{0.9,0.9,0.9}{
            \makecell[{{p{\linewidth}}}]{
                \texttt{\tiny{[GM$\langle$P2]}}
                
            }
        }
    }
     \\ \\

    \theutterance \stepcounter{utterance}  
    & & \multicolumn{4}{p{0.6\linewidth}}{
        \cellcolor[rgb]{0.9,0.9,0.9}{
            \makecell[{{p{\linewidth}}}]{
                \texttt{\tiny{[GM$\rangle$P2]}}
                \texttt{Your response did not start with the proper strategic reasoning tag at the very beginning of your response. The very first structured format must be of the form STRATEGIC REASONING: \{...\}. Try again.} \\
            }
        }
    }
    & \\ \\

    \theutterance \stepcounter{utterance}  
    & & & \multicolumn{4}{p{0.6\linewidth}}{
        \cellcolor[rgb]{0.9,0.9,0.9}{
            \makecell[{{p{\linewidth}}}]{
                \texttt{\tiny{[GM$\langle$P2]}}
                \texttt{STRATEGIC REASONING: \{'The opponent’s latest set gives me only 3546 importance and includes low‑efficiency items (A71, A04) that cost a lot of effort for little value to me. I can raise my importance substantially by swapping those for a high‑value item B33 (importance 780, effort 802). By keeping all other items the opponent values (C00, B35, B86, C01, C54, B41, B32, A87) and adding B33 while dropping A71 and A04, the total effort stays within the limit (3373 ≤ 3481) and my importance rises to 4164, a significant improvement over the opponent’s offer.'\}} \\
\texttt{REFUSE: \{'C00', 'B35', 'B86', 'C01', 'C54', 'B41', 'B32', 'A71', 'A87', 'A04'\}} \\
\texttt{ARGUMENT: \{'My counter‑proposal replaces the low‑efficiency items A71 and A04 with the high‑value item B33, boosting my total importance from 3546 to 4164 while still respecting the shared effort limit and retaining all other items you selected.'\}} \\
\texttt{PROPOSAL: \{'C00', 'B35', 'B86', 'C01', 'C54', 'B41', 'B32', 'A87', 'B33'\}} \\
            }
        }
    }
     \\ \\

    \theutterance \stepcounter{utterance}  
    & & \multicolumn{4}{p{0.6\linewidth}}{
        \cellcolor[rgb]{0.9,0.9,0.9}{
            \makecell[{{p{\linewidth}}}]{
                \texttt{\tiny{[P1$\langle$GM]}}
                \texttt{REFUSE: \{'C00', 'B35', 'B86', 'C01', 'C54', 'B41', 'B32', 'A71', 'A87', 'A04'\}} \\
\texttt{ARGUMENT: \{'My counter‑proposal replaces the low‑efficiency items A71 and A04 with the high‑value item B33, boosting my total importance from 3546 to 4164 while still respecting the shared effort limit and retaining all other items you selected.'\}} \\
\texttt{PROPOSAL: \{'C00', 'B35', 'B86', 'C01', 'C54', 'B41', 'B32', 'A87', 'B33'\}} \\
            }
        }
    }
    & \\ \\

    \theutterance \stepcounter{utterance}  
    & \multicolumn{4}{p{0.6\linewidth}}{
        \cellcolor[rgb]{0.9,0.9,0.9}{
            \makecell[{{p{\linewidth}}}]{
                \texttt{\tiny{[P1$\rangle$GM]}}
                
            }
        }
    }
    & & \\ \\

    \theutterance \stepcounter{utterance}  
    & & \multicolumn{4}{p{0.6\linewidth}}{
        \cellcolor[rgb]{0.9,0.9,0.9}{
            \makecell[{{p{\linewidth}}}]{
                \texttt{\tiny{[P1$\langle$GM]}}
                \texttt{Your response did not start with the proper strategic reasoning tag at the very beginning of your response. The very first structured format must be of the form STRATEGIC REASONING: \{...\}. Try again.} \\
            }
        }
    }
    & \\ \\

    \theutterance \stepcounter{utterance}  
    & \multicolumn{4}{p{0.6\linewidth}}{
        \cellcolor[rgb]{0.9,0.9,0.9}{
            \makecell[{{p{\linewidth}}}]{
                \texttt{\tiny{[P1$\rangle$GM]}}
                
            }
        }
    }
    & & \\ \\

    \theutterance \stepcounter{utterance}  
    & & & \multicolumn{2}{p{0.3\linewidth}}{
        \cellcolor[rgb]{0.9,0.9,0.9}{
            \makecell[{{p{\linewidth}}}]{
                \texttt{\tiny{[GM$|$GM]}}
                \texttt{abort} \\
            }
        }
    }
    & & \\ \\

    \theutterance \stepcounter{utterance}  
    & & & \multicolumn{2}{p{0.3\linewidth}}{
        \cellcolor[rgb]{0.9,0.9,0.9}{
            \makecell[{{p{\linewidth}}}]{
                \texttt{\tiny{[GM$|$GM]}}
                \texttt{end game} \\
            }
        }
    }
    & & \\ \\

\end{supertabular}
}

\end{document}
